We have seen how we can take on the integration problem for formal libraries by representing them in a unifying meta-framework (\autoref{ch:import}), using a flexible logical framework to specify their foundations (\autoref{ch:lf}), aligning their contents (\autoref{sec:crosslib:alignments}) and implementing foundation-independent alignment-aware algorithms for identifying overlap (\autoref{sec:crosslib:vf}) suggestive of new alignments, translating expressions across libraries (\autoref{sec:crosslib:translate}) and modularize library content by refactoring along theory intersections (\autoref{sec:crosslib:intersect}).

The knowledge management techniques described in \autoref{ch:crosslib} are certainly not exhaustive; many more services are imaginable, but the ones presented in this thesis demonstrate the effectiveness of the general approach to library integration described herein. Furthermore, these facilitate the aims of the OAF (see \autoref{sec:intro:objectives}) and \ODK (see \autoref{sec:intro:prelim:odk} ) projects, as demonstrated in \autoref{sec:crosslib:translate:ex}. Additionally they represent steps towards a system integrating all aspects of the tetrapod (see \autoref{ch:intro}); demonstrably \emph{inference}, \emph{computation}, \emph{tabulation} and \emph{organization}.

\paragraph{} Naturally, this opens up many avenues for future research and improvements, which fall broadly into three categories:

\paragraph{1. Fundamentals} Many implementations described in this thesis are prototypical and can be optimized or otherwise improved. In particular:
\begin{enumerate}
	\item The implementation of alignments (see \autoref{sec:crosslib:alignments}) is not deeply integrated into the \mmt system and offers no change management facilities. As mentioned in \autoref{sec:crosslib:alignments:manually}, many of the alignments we once had are by now outdated for various reasons, suggesting that additional services are needed that e.g. keep track of (and adequately adapt) changes in URIs, check for the existence (or maybe even adequacy) of aligned symbols, allow for globally replacing URIs in multiple alignments, etc.
	\item The current syntax for alignments is, while simple and convenient, not exhaustive for many types of alignments, such as those that need simple arithmetic operations on arguments (as the dihedral groups described in \autoref{sec:crosslib:translate:ex}), or the cases where a function $f(x)$ is aligned with a function $g(t,x)$ for a fixed expression $t$. Currently, this introduces a need for programmatic translators, but a more expressive syntax could obviate that need in many situations.
	\item The implementation of alignment-based translation presented in \autoref{sec:crosslib:translate} is naive and potentially inefficient. While this does not show for the use cases that have occured in practie so far, I strongly suspect that this might become inconvenient once the number of available alignments and other translators grows, or the input expressions reach a certain degree of complexity.
	\item I mentioned in the introduction to \autoref{ch:lf} that it is desirable to not have inference rules in a logical framework active that are not needed for some specific content -- one of the reasons why we prefer a modular approach to our logical frameworks. This is also the reason we shied away from a naive implementation of declared subtyping in \autoref{sec:lf:lfx:subtp}. To expand on that claim, the reason is that naively applying inference rules can be computationally expensive and slow down the solver. This is not just the case for unnecessary inference rules, but also for more basic conflicting strategies for checking a judgement (\emph{critical pairs}). A simple example is the basic approach of a bi-directional type checking algorithm:
	
	To check a judgement $\hastype tT$ we have two strategies: We can either simplify $T$ until we find an applicable foundation-dependent typing rule specifically for (the head symbol of) $T$, or we can try to infer the principal type $T'$ of $t$ and check $\subtp {T'}T$. In most cases, the existence of a typing rule implies that the former strategy will quickly yield a result, whereas inferring the type can be expensive, and hence the former strategy is attempted first. However, imagine $T$ being a single reference to a constant with a complex definiens, and $t$ being a function application $f(a)$, where the return type of $f$ is $T$. In this case, inferring the principal type of $t$ is fast and the subsequent check $\subtp TT$ trivial, whereas simplifying $T$ leads to definition expansion and subsequent further simplifications, which will ultimately need to be repeated for the inferred type of $t$ as well, since a function application can only ever check against any type via the inference strategy. This kind of situation seems to occur regularly in high-level formalizations, and is correspondingly often encountered in the development of the Math-in-the-Middle library.
	
	To remedy this problem, I am currently actively working on a reimplementation of the \mmt solver algorithm using multithreading; not just to exploit multicore processors, but to be able to compute conflicting checking strategies on separate threads in parallel, such that if \emph{either} strategy terminates quickly, the checking algorithm can return a result just as quickly. This would obviate the current advantage of being cautious with adding potentially slow inference rules.
\end{enumerate}

\paragraph{2. Scaling Up} Many of the implementations and approaches presented in this thesis scale in usefulness and coverage depending on a certain amount of available data, features and implementations. All of these can and should be scaled up in the future, such as:
\begin{enumerate}
	\item extending our logical framework by more desirable features (\autoref{ch:lf}),
	\item Importing more formal libraries and system ontologies into \ommt (\autoref{ch:import}),
	\item expanding the Math-in-the-Middle library by more content and interface theories (\autoref{sec:crosslib:translate}),
	\item collecting more alignments between libraries (\autoref{sec:crosslib:alignments}),
	\item covering more library contents with dedicated translators where needed (\autoref{sec:crosslib:translate}),
	\item investigating both the available and additional libraries for more useful preprocessing strategies to increase the coverage and usefulness of the view finder algorithm (\autoref{sec:crosslib:vf}),
\end{enumerate}

\paragraph{3. Workflows} One of the major obstacles in using all of the techniques in this thesis is a certain lack of tool support with respect to specific workflows for various use cases. Especially the algorithms presented in \autoref{ch:crosslib} are primarily usable by programmatically calling methods in the \mmt API, or via commands in the \mmt shell. Exemplary, in \autoref{sec:crosslib:vf} my coauthors and I (on the original publication) already conjecture on possible applications for the view finder algorithm, most of which are predicated on the existence of a bespoke user interface for that particular use case.
Designing such interfaces entails thinking deeply about realistic and convenient workflows for different situations and users. 

This also applies to \mmt's surface syntax, which requires a very reasonable, but (for example, for working mathematicians not familiar with formal systems) potentially prohibitive amount of time and effort to get accustomed to.

I have written an \mmt plugin for IntelliJ IDEA\footnote{\url{https://www.jetbrains.com/idea/}, the plugin itself is available at \url{https://github.com/UniFormal/IntelliJ-MMT}}, a java-based IDE which offers a rich plugin API, to experiment with integrating various generic services in a more flexible IDE than the only previously supported one, jEdit. In the future, I would like to investigate the possibility of using e.g. s\TeX~\cite{stex} or possibly even plain \LaTeX\ as an alternative input language for formal \mmt content, as well as designing a unifying online IDE on MathHub for working with \mmt and its connected systems and libraries, obviating the need to set up and use \mmt (and the multitude of systems with which \mmt can interact) locally -- a process which (even using docker images, jupyter notebooks and similar solutions) can be prohibitively complicated for non-experts, especially when lacking a unifying user interface.